\section{Method}
\label{sec:method}

Let $f_T:\mathcal{X}\rightarrow\mathbb{R}^{d_T}$ be a frozen teacher and
$f_S(\cdot;\theta):\mathcal{X}\rightarrow\mathbb{R}^{d_S}$ a student with
parameters $\theta$ and $d_T>d_S$. We represent embeddings as row vectors
throughout. We run the teacher once over an unlabeled corpus
$\{x_i\}_{i=1}^{N}$ and cache its normalized final sentence embeddings
\begin{equation}
    \mathbf{t}_i=\operatorname{norm}\!\left(f_T(x_i)\right)
    \in\mathbb{S}^{d_T-1},
\end{equation}
where $\operatorname{norm}(\mathbf{x})=\mathbf{x}/\lVert\mathbf{x}\rVert_2$.
The corresponding student embedding is
$\mathbf{z}_i(\theta)=\operatorname{norm}(f_S(x_i;\theta))\in
\mathbb{S}^{d_S-1}$; we suppress $\theta$ when it is clear from context.
GATE-KD primarily transfers teacher knowledge through an
extrinsic interface that expresses each teacher endpoint in coordinates the
student can match pointwise. It augments this interface with an intrinsic
regularizer that requires neither shared coordinates nor equal ambient
dimension.

\subsection{Extrinsic Endpoint Interface}
\label{sec:method-targets}

\paragraph{Subspace selection.}
Let $\mu=N^{-1}\sum_i\mathbf{t}_i$. We fit the leading $d_S$ principal
directions of the centered cache and collect them as
$P_{\mathrm{PCA}}\in\mathbb{R}^{d_T\times d_S}$, with
$P_{\mathrm{PCA}}^\top P_{\mathrm{PCA}}=I_{d_S}$. Centering is used to
estimate the principal directions; in the evaluated recipe, the cached vectors
are not centered when the projection is applied. The normalized coordinates
before gauge alignment are
\begin{equation}
    \mathbf{y}_i
    =\operatorname{norm}\!\left(\mathbf{t}_iP_{\mathrm{PCA}}\right)
    \in\mathbb{S}^{d_S-1}.
    \label{eq:pca-target}
\end{equation}
The PCA subspace is fitted once on the full teacher cache and remains fixed
throughout training. Centering the cache to estimate covariance directions and
subtracting the mean when applying the interface are distinct operations: the
former selects a subspace, whereas the latter would make the target transform
affine and alter its angular geometry after row normalization. We use only the
former, preserving a linear endpoint interface. Appendix~\ref{app:pca-controls}
evaluates both choices together with whitening and random-subspace controls.

\paragraph{Gauge alignment.}
The basis used to express the retained subspace affects a pointwise loss even
though it does not affect the geometry within that subspace. We resolve this
coordinate ambiguity with orthogonal Procrustes alignment
\citep{schonemann1966procrustes,bhattarai2025shortcut}. We treat the orientation
as a nuisance coordinate variable rather than as teacher content. A gauge-fit
set $\mathcal{C}$ is selected before training and reused unchanged; in the main
experiments it is the complete distillation corpus. Let
$Y_{\mathcal{C}}$ contain $\mathbf{y}_i$ as rows and let
$Z_{\mathcal{C}}^{(e)}$ contain the corresponding student embeddings
$\mathbf{z}_i^{(e)}$ at the start of epoch $e$. We compute
\begin{equation}
    \begin{aligned}
    U\Sigma V^\top
      &=\operatorname{svd}\!\left(Y_{\mathcal{C}}^\top
        Z_{\mathcal{C}}^{(e)}\right),\\
    R^{(e)}=UV^\top
      &\in\underset{R\in O(d_S)}{\arg\min}\,
        \left\lVert Y_{\mathcal{C}}R-Z_{\mathcal{C}}^{(e)}\right\rVert_F^2.
    \end{aligned}
    \label{eq:gauge-alignment}
\end{equation}
The initialized student determines $R^{(0)}$. After each non-final epoch, we
re-encode the same gauge-fit sentences and refit the gauge for the following
epoch. The target used within epoch $e$ is therefore
\begin{equation}
    \boldsymbol{\tau}_i^{(e)}=\mathbf{y}_iR^{(e)}.
    \label{eq:target}
\end{equation}
Each target is fixed within an epoch. The refit is a closed-form coordinate
update rather than a learned projector. The minimizer in
Equation~\ref{eq:gauge-alignment} is unique exactly when the cross-covariance
$Y_{\mathcal{C}}^\top Z_{\mathcal{C}}^{(e)}$ has full rank
\citep{schonemann1966procrustes}; in directions with vanishing singular
values, $UV^\top$ is one of several minimizers and the returned orientation
depends on the SVD convention.

\paragraph{Relation to the endpoint objective.}
Because the rows of $Y_{\mathcal{C}}$ and $Z_{\mathcal{C}}^{(e)}$ have unit
norm, their cosine endpoint objective satisfies
\begin{equation}
    \frac{1}{|\mathcal{C}|}\sum_{i\in\mathcal{C}}
    \left(1-\left\langle\mathbf{z}_i^{(e)},\mathbf{y}_iR\right\rangle\right)
    =\frac{1}{2|\mathcal{C}|}
      \left\lVert Z_{\mathcal{C}}^{(e)}-Y_{\mathcal{C}}R\right\rVert_F^2.
    \label{eq:endpoint-procrustes}
\end{equation}
\begin{proposition}[Optimality and invariance of the fitted gauge]
\label{prop:gauge-optimality}
For row-normalized $Y_{\mathcal{C}}$ and $Z_{\mathcal{C}}^{(e)}$, the update
$R^{(e)}$ in Equation~\ref{eq:gauge-alignment} globally minimizes the gauge-fit
cosine endpoint loss over $O(d_S)$. Moreover, for every projected target matrix
$Y$ and every $R\in O(d_S)$,
\begin{equation}
    (Y R)(Y R)^\top=YY^\top.
    \label{eq:gauge-invariance}
\end{equation}
Thus fitting the gauge changes the coordinate presentation but not sample
correspondence, pairwise cosines, chord distances, or persistent-homology
signatures of the projected teacher targets.
\end{proposition}
The proof is given in Appendix~\ref{app:supporting-proofs}.

PCA supplies a fixed student-width subspace, while Procrustes determines how
that fixed geometry is presented in the student's coordinates. The update does
not select a new subspace or adapt the teacher's within-target geometry. At each
refit boundary, with the student fixed, it cannot increase the gauge-fit
endpoint loss and leaves the native-space $H_0$ term unchanged. This conditional
guarantee does not imply monotonic epoch-to-epoch training because the student
changes between refits.

\subsection{Native-Space \texorpdfstring{$H_0$}{H0} Persistence-Summary Regularization}
\label{sec:method-topology}

The endpoint interface can reproduce only the geometry retained after
projection. Although $P_{\mathrm{PCA}}$ has orthonormal columns,
$P_{\mathrm{PCA}}P_{\mathrm{PCA}}^\top$ is a rank-$d_S$ projector, so the
cosine Gram matrix of the normalized coordinates in
Equation~\ref{eq:pca-target} need not equal that of the original teacher
embeddings. This loss of geometry is unavoidable whenever the normalized
teacher cache has rank above the student width.
\begin{proposition}[Rank-relaxed lower bound on Gram distortion]
\label{prop:gram-lower-bound}
Let $T\in\mathbb{R}^{N\times d_T}$ collect the normalized teacher cache
embeddings, with singular values
$\sigma_1(T)\geq\cdots\geq\sigma_r(T)>0$. For every $Z\in
\mathbb{R}^{N\times d_S}$ collecting student embeddings of the same $N$
sentences,
\begin{equation}
    \left\lVert TT^\top-ZZ^\top\right\rVert_F^2
    \geq \sum_{j>d_S}\sigma_j(T)^4.
    \label{eq:gram-lower-bound}
\end{equation}
Thus exact preservation of the teacher cosine Gram matrix is impossible when
$\operatorname{rank}(T)>d_S$. The bound relaxes the unit-diagonal constraint on
$ZZ^\top$ and therefore need not be tight for normalized student embeddings.
\end{proposition}
The proof is given in Appendix~\ref{app:supporting-proofs}.

Rather than attempting to reconstruct the full teacher Gram geometry, which is
generally impossible at student rank, we match a lower-complexity,
coordinate-free summary of its connectivity scales. We therefore add a
persistence-summary regularizer defined directly in the teacher's native space.

For a minibatch $\mathcal{B}\subseteq\{1,\ldots,N\}$ of size
$|\mathcal{B}|=B$, we use the chord distance
\begin{equation}
    d(\mathbf{u},\mathbf{v})
    =\sqrt{2-2\langle\mathbf{u},\mathbf{v}\rangle}.
\end{equation}
We parameterize the Vietoris--Rips filtration by the distance threshold: an
edge enters at scale $a$ when its distance is at most $a$.
The $B-1$ finite death times of degree-zero Vietoris--Rips persistence are the
edge weights of the point cloud's minimum spanning tree
\citep{hofer2020topologically}. Let
$q^S_{(1)}\leq\cdots\leq q^S_{(B-1)}$ and
$q^T_{(1)}\leq\cdots\leq q^T_{(B-1)}$ denote the sorted death times
of the student embeddings $\{\mathbf{z}_i:i\in\mathcal{B}\}$ and the original
teacher embeddings $\{\mathbf{t}_i:i\in\mathcal{B}\}$. We define
\begin{equation}
    \mathcal{L}_{H_0}
    =\frac{1}{B-1}\sum_{j=1}^{B-1}
      \left(q^S_{(j)}-q^T_{(j)}\right)^2.
    \label{eq:topological-loss}
\end{equation}
For a precise interpretation, define the empirical distributions of finite
connectivity scales
\begin{equation}
    \nu_S=\frac{1}{B-1}\sum_{j=1}^{B-1}\delta_{q^S_{(j)}},
    \qquad
    \nu_T=\frac{1}{B-1}\sum_{j=1}^{B-1}\delta_{q^T_{(j)}},
\end{equation}
where $\delta_a$ denotes a Dirac mass at $a$.
\begin{proposition}[Transport interpretation and stability]
\label{prop:h0-transport-stability}
The $H_0$ loss equals the squared one-dimensional Wasserstein distance between
the empirical death-scale measures,
\begin{equation}
    \mathcal{L}_{H_0}=W_2^2(\nu_S,\nu_T).
    \label{eq:topological-wasserstein}
\end{equation}
Moreover, if the two distance matrices on the corresponding batch satisfy
\begin{equation}
    \epsilon=\max_{i,j\in\mathcal{B}}
    \left|d(\mathbf{z}_i,\mathbf{z}_j)
          -d(\mathbf{t}_i,\mathbf{t}_j)\right|,
\end{equation}
then $\mathcal{L}_{H_0}\leq\epsilon^2$.
\end{proposition}
The proof is given in Appendix~\ref{app:supporting-proofs}.

The teacher diagram is computed from $\mathbf{t}_i\in\mathbb{R}^{d_T}$ rather
than $\mathbf{y}_i$ or $\boldsymbol{\tau}_i^{(e)}$. The loss therefore needs
neither a projection nor a shared orientation. It is differentiable almost
everywhere away from distance ties and changes in the active minimum spanning
tree \citep{hofer2019connectivity}. In implementation, the discrete tree
selection and sorting permutation are treated as fixed during backpropagation,
while gradients flow through the selected student edge weights.
The stability result is one-way. Equation~\ref{eq:topological-loss} matches the
multiset of $H_0$ death times but does not impose correspondence between
individual minimum-spanning-tree edges, recover the full metric geometry, or
preserve higher-dimensional topological features.

Algorithm~\ref{alg:h0-loss} makes the nonsmooth implementation explicit. The
teacher branch is evaluated without gradients. On the student branch, the MST
edge selection and sorting permutation are detached, while the selected edge
weights remain differentiable. Exact distance ties are passed to SciPy in
row-major upper-triangular order. Although the selected active tree can depend
on this library tie-break, every valid MST has the same sorted multiset of edge
weights; the backward pass follows the particular active tree returned in that
forward pass.

\begin{algorithm}[t]
\caption{Native-space $H_0$ connectivity-scale loss for one minibatch}
\label{alg:h0-loss}
\begin{algorithmic}[1]
\Require Row-normalized teacher and student batches $T_{\mathcal B}$ and
$Z_{\mathcal B}$; numerical floor $\varepsilon$.
\For{$A\in\{T_{\mathcal B},Z_{\mathcal B}\}$}
    \State Compute $G\gets AA^\top$ in FP32 and symmetrize
    $G\gets(G+G^\top)/2$.
    \State Compute $D_{ij}\gets\sqrt{\max(2-2G_{ij},\varepsilon)}$ and set
    $D_{ii}\gets0$.
    \State On a detached CPU float64 copy, add the same positive constant to
    every off-diagonal edge and select an undirected MST with SciPy from the
    upper triangle.
    \State Gather the selected weights from the unshifted $D$ and sort them to
    obtain $q^A_{(1:B-1)}$.
\EndFor
\State \Return $(B-1)^{-1}\sum_{j=1}^{B-1}(q^Z_{(j)}-q^T_{(j)})^2$.
\end{algorithmic}
\end{algorithm}

\subsection{Objective and Alternating Optimization}
\label{sec:method-objective}

The endpoint loss for the same training minibatch $\mathcal{B}$ is
\begin{equation}
    \mathcal{L}_{\mathrm{end}}^{(e)}
    =\frac{1}{B}\sum_{i\in\mathcal{B}}
      \left(1-\left\langle
      \mathbf{z}_i,\boldsymbol{\tau}_i^{(e)}\right\rangle\right).
    \label{eq:endpoint-loss}
\end{equation}
GATE-KD optimizes the two-term objective
\begin{equation}
    \mathcal{L}^{(e)}
    =\mathcal{L}_{\mathrm{end}}^{(e)}
    +\lambda_{\mathrm{topo}}\mathcal{L}_{H_0}.
    \label{eq:objective}
\end{equation}
Here, $\mathcal{L}_{H_0}$ is evaluated on the same minibatch using the student
embeddings and the corresponding unprojected teacher embeddings.

Optimization alternates between stochastic updates of the student parameters
and exact updates of the nuisance coordinate gauge. It can be viewed as an
inexact block-coordinate procedure over $\theta$ and $R\in O(d_S)$: with the
student fixed, Equation~\ref{eq:gauge-alignment} solves the endpoint term
exactly over $R$ on $\mathcal{C}$; with the gauge fixed, stochastic minibatch
updates modify $\theta$. During epoch $e$, $R^{(e)}$ is held
fixed while the student is optimized with Equation~\ref{eq:objective}. After
epoch $e$, unless it is the final epoch, the updated student re-encodes
$\mathcal{C}$ and Equation~\ref{eq:gauge-alignment} produces $R^{(e+1)}$ for
the next epoch. The $H_0$ term is unaffected by these gauge updates because it
is computed from the native teacher embeddings.

Teacher embeddings and the PCA subspace are computed once before student
optimization. Ordinary training steps use the cached teacher embeddings and
require no online teacher forward pass. Each gauge update adds one student pass
over the gauge-fit set and an SVD of a $d_S\times d_S$ matrix. This joint view
does not by itself imply that repeated refitting improves generalization; the
controlled analysis in Section~\ref{sec:analysis-interface} treats it as a
secondary refinement beyond the initially fitted gauge.
